\documentclass[a4paper,12pt]{article}
\usepackage[utf8]{inputenc}
\usepackage[T1]{fontenc}
\usepackage[portuguese]{babel}
\usepackage{geometry}
\geometry{margin=2.5cm}
\usepackage{enumitem}
\usepackage{xcolor}
\usepackage{tcolorbox}

\definecolor{ifscgreen}{HTML}{329632}

\begin{document}

\begin{minipage}{0.2\textwidth}
    % Espaço para logo se necessário
\end{minipage}
\begin{minipage}{0.8\textwidth}
    \begin{flushright}
        \textbf{IFSC - Campus São José} \\
        \textbf{Superior em Sistemas para Internet} \\
        \textbf{UC: Programação para Internet} \\
        \textbf{Prof. Andre Chameo}
    \end{flushright}
\end{minipage}

\vspace{1cm}

\begin{center}
    \begin{tcolorbox}[colback=ifscgreen!10,colframe=ifscgreen,title=\centering\large\bfseries ATIVIDADE PRÁTICA: DESENVOLVIMENTO WEB FRONT-END]
        \centering\small Trabalho em Duplas - Prazo: [Definir Data]
    \end{tcolorbox}
\end{center}

\section{Objetivo}
O objetivo desta atividade é aplicar de forma integrada os conhecimentos de \textbf{HTML5 Semântico}, \textbf{CSS3 (Flexbox/Grid/Responsividade)} e \textbf{JavaScript (DOM/Eventos)} para construir uma interface web moderna, funcional e esteticamente atraente.

\section{Orientações Gerais}
\begin{itemize}
    \item A atividade deve ser realizada em \textbf{duplas}.
    \item O projeto deve ser versionado no \textbf{GitHub} (um repositório por dupla).
    \item É obrigatório o uso de design responsivo (Mobile First).
    \item Não é permitido o uso de frameworks CSS (Bootstrap, Tailwind) ou bibliotecas externas, exceto fontes (Google Fonts) e ícones (FontAwesome/SVG).
\end{itemize}

\section{Temáticas Sugeridas}
Escolha uma das temáticas abaixo para o desenvolvimento do projeto:

\begin{enumerate}[label=\textbf{Tema \arabic*:}]
    \item \textbf{E-Commerce de Nicho:} Uma loja virtual para um produto específico (ex: plantas, cafés, vinis). Deve conter cabeçalho, galeria de produtos com filtros e um modal de "carrinho".
    \item \textbf{Dashboard de Monitoramento:} Um painel administrativo para um sistema de Smart Home ou Saúde. Deve conter menu lateral, cartões de métricas (temperatura, batimentos) e tabelas de histórico.
    \item \textbf{Landing Page de Evento Tecnológico:} Uma página de inscrição para um congresso ou hackathon. Deve conter hero section, cronograma, palestrantes e formulário de inscrição com validação.
    \item \textbf{Portal de Conteúdo/Notícias:} Um site de notícias focado em um tema (ex: Sustentabilidade, Games). Deve focar em hierarquia tipográfica, grid de notícias e barra lateral de "mais lidas".
\end{enumerate}

\section{Requisitos Técnicos}
\begin{itemize}
    \item \textbf{HTML:} Uso correto de tags semânticas (\texttt{main}, \texttt{section}, \texttt{article}, \texttt{aside}, \texttt{nav}).
    \item \textbf{CSS:} Layout construído com Flexbox e/ou Grid Layout. Uso de variáveis CSS e pseudo-classes (\texttt{:hover}, \texttt{:active}).
    \item \textbf{JavaScript:} Implementação de ao menos uma funcionalidade interativa (ex: carrossel, validação de form, alternância de tema dark/light ou filtragem dinâmica).
\end{itemize}

\section{Critérios de Avaliação}
\begin{enumerate}
    \item \textbf{Estrutura e Semântica (20\%):} Organização do HTML.
    \item \textbf{Design e Estética (30\%):} Harmonia de cores, tipografia e acabamento visual.
    \item \textbf{Responsividade (25\%):} Adaptação para diferentes tamanhos de tela.
    \item \textbf{Funcionalidade JS (25\%):} Lógica e interação com o usuário via DOM.
\end{enumerate}

\end{document}
